\documentclass[11pt]{article}

% Scalable Type 1 fonts (Latin Modern): required for microtype's font
% expansion on TeX installs without cm-super, and keeps the PDF text
% vector (crisp in browser viewers) rather than bitmap Type 3.
\usepackage[T1]{fontenc}
\usepackage{lmodern}
\usepackage[a4paper,margin=1in]{geometry}
\usepackage{amsmath,amssymb,amsthm}
\usepackage{mathtools}
\usepackage{microtype}
\usepackage{booktabs}
\usepackage{array}
\usepackage{graphicx}
\usepackage{tikz}
\usepackage{pgfplots}
\pgfplotsset{compat=1.18}
\usetikzlibrary{arrows.meta,positioning,calc,shapes.geometric,shapes.misc,fit,backgrounds,decorations.pathreplacing,patterns}
\usepackage[hidelinks]{hyperref}
\usepackage{xcolor}
\usepackage{enumitem}
\usepackage{caption}
\usepackage{tabularx}
\usepackage{url}
\usepackage{algorithm}
\usepackage{algpseudocode}

\definecolor{accent}{HTML}{1f3a93}
\definecolor{accent2}{HTML}{c0392b}
\definecolor{accent3}{HTML}{27ae60}
\definecolor{rule}{HTML}{888888}
\definecolor{lightgray}{HTML}{dcdcdc}

\hypersetup{colorlinks=true, linkcolor=accent, citecolor=accent, urlcolor=accent}
\emergencystretch=3em
\sloppy

\theoremstyle{plain}
\newtheorem{theorem}{Theorem}
\newtheorem{proposition}[theorem]{Proposition}
\newtheorem{lemma}[theorem]{Lemma}
\theoremstyle{definition}
\newtheorem{definition}[theorem]{Definition}
\newtheorem{remark}[theorem]{Remark}

\newcommand{\bpp}{\mathrm{bpp}}
\newcommand{\kl}{D_{\mathrm{KL}}}
\newcommand{\Linears}{\mathcal{L}}
\newcommand{\Formats}{\mathcal{F}}
\newcommand{\Assignment}{\mathcal{A}}
\newcommand{\R}{\mathbb{R}}
\newcommand{\E}{\mathbb{E}}
\newcommand{\aura}{\textsc{Aura}}
\newcommand{\dW}{\Delta W}

% =====================================================================
% CANONICAL paper as of 2026-06-05 (PrismaSCOUT retired; old PrismaSCOUT paper
% preserved under paper/archive/prismascout_paper_2026-06-05.tex).
% DRAFT v0.1. Authored by Claude in deliberation with Codex
% (high-reasoning pass: .claude/codex-aura-paper-2026-06-05.md). Spine:
% AURA leads as the method; the bf16-floor reframe is a softened, evidence-
% labelled "measurement-artifact study"; weight-only and bf16-dW limitations
% owned up front. ALL numbers trace to /home/rob/dq-runs/{xlayer_rung0,
% aura-27b-flagship,aura-sat-validate-4b,aura-blocker-runs} and the auto-memory;
% the 2026-06-12 codex review re-derived them from raw logs
% (docs/archive/codex-process/CODEX_QUEUE_FINDINGS.md ~L909). xlayer_rung0/ was deleted in a July disk cleanup AFTER
% that verification. 2026-07-13 accuracy pass: baseline re-labelled
% "then-production-default" (M6 flipped the default to h_trace*weight_mse,
% commit 8588036), ToolEvalBench behavioral corroboration added, MoE
% route-flip boundary added to limitations, bib TODOs resolved from arXiv.
% =====================================================================

\title{\bfseries \textsc{Aura}: Production-Faithful KL--Fisher Allocation\\
       for Mixed-Precision Large Language Model Quantization}

\author{Robert Tand\\
        \textit{Independent Researcher}\\
        \texttt{robert.tand@icloud.com}}

\date{July 2026 \quad\textsc{(draft)}}

\begin{document}
\maketitle

\begin{abstract}
Mixed-precision quantization assigns each linear layer of a transformer its own
serving format under a global bit budget. The dominant view in the allocation
literature is that the per-Linear cost surrogates which drive this assignment are
\emph{structurally} biased --- that quantization error is non-additive across
layers --- and that the remedy is to \emph{model} the cross-layer interaction
(pairwise integer-quadratic programming, second-order ILP, cooperative-game
allocation). We report a different, and we argue simpler, account for the
serving regimes that matter on current hardware ($\ge 4$ bits). We introduce
\textsc{Aura}, a per-Linear allocation cost equal to an $O(N)$ randomized
estimator of the downstream Kullback--Leibler (KL) contribution of quantizing
that Linear to a given format, evaluated against the production-rendered weight
error the exporter ships (the GPTQ+JSO bytes, not an idealized round; stored in
bf16 in the current implementation). \textsc{Aura} is additive across Linears,
solved as a multiple-choice knapsack, and --- critically --- every candidate it
proposes is gated by a real, end-to-end KL measurement on a held-out split.
Against the then-production-default Fisher-trace$\times$output-MSE baseline,
\textsc{Aura} lowers served vLLM KL-vs-BF16 by $\mathbf{17.9\%}$ at matched
$5.99$~bpp on Qwen3.6-27B (replicated on two disjoint held-out corpora); against a
weaker output-MSE baseline the $4.85$~bpp knee on Qwen3-4B improves
$\mathbf{38\text{--}40\%}$. The 27B gap is larger at $5.99$ than at $5.25$~bpp on a
single seed (confidence intervals pending). We are explicit
that this is a distributional-fidelity (KL) win, not a raw-perplexity win: at
$6$~bpp perplexity is near-saturated for both arms ($+2.3\%$ vs.\ $+2.7\%$ over
BF16). A supporting measurement study explains \emph{why} an additive cost
suffices: much of the apparent additive overshoot reported in the cross-layer
literature is a \emph{numerical-precision floor} of bf16 KL differencing --- in
fp32 (identical to fp64 to five significant figures) the per-Linear unary KLs
are nearly additive, with a small but nonzero genuine residual that the
real-KL selector absorbs. Finally, we give the operational rate--distortion
geometry: the achievable frontier is the integral of the sorted per-Linear
cost-per-bit, so it is convex, and a \emph{sharp} knee is a property of the
format \emph{menu} (an intermediate-efficiency rung such as FP8), not of the
sensitivity distribution --- which is why a two-rung NVFP4/BF16 frontier looks
linear on a log-KL axis even though a better cost still wins by shifting it
\emph{down}. We position \textsc{Aura} as a production-faithful realization of
the well-studied Fisher/KL sensitivity family, not as new sensitivity theory,
and we report the limitations (weight-only scoring under W4A4 activation
quantization; bf16-stored rendering error) that a faithful account requires.
In doing so we retire the multi-level cost cascade of our prior
\textsc{PrismaSCOUT} system in favour of a single additive cost plus served-KL
selection.
\end{abstract}

% =====================================================================
\section{Introduction}\label{sec:intro}

Post-training quantization of large language models splits cleanly into two
orthogonal questions. The \emph{local} question --- given a fixed format, how do
you round one Linear best? --- is well studied: GPTQ \cite{gptq2022},
AutoRound \cite{autoround2023}, rotation preprocessors, and most recently
ParoQuant \cite{paroquant}. The \emph{global} question --- how many bits should
each Linear get, and in which hardware format? --- is allocation, and is the
subject of this paper. A heterogeneous per-Linear assignment over a menu such as
$\{$BF16, FP8, NVFP4$\}$ extracts quality no single-format method can, because
sensitivity is wildly unequal across a transformer's Linears.

The allocation literature is unified by one premise: per-Linear sensitivity
surrogates (diagonal Hessian traces, Fisher MSEs, output MSEs) are biased
estimators of joint quantization error, and the bias is \emph{cross-layer}. The
standard remedy is to model it: the cross-layer-dependency integer-quadratic
program of CLADO \cite{clado2023}, the second-order ILP of HAWQ-V3
\cite{hawqv3}, the Shapley cooperative-game allocator of CoopQ \cite{coopq}.
A prior system of ours, \textsc{PrismaQuant}/\textsc{PrismaSCOUT}
\cite{prismaquant}, took the position that one need not model the interaction at
all if every shippable candidate is re-scored end-to-end --- ``surrogates
generate, real measurement selects'' --- but still paid for a three-level cost
cascade ($L_1$ probe, $L_2$ perturbed-activation fixed point, $L_3$ propagated
end-KL) whose purpose was to chase the same non-additivity. \textbf{This paper
retires that cascade.} We keep its real-KL \emph{selector} --- the one component
that observes the true objective --- and replace the $L_2/L_3$ \emph{modeling}
with a single additive cost, on the evidence (Section~\ref{sec:additivity}) that
at $\ge 4$ bits the non-additivity it modeled is largely a measurement-precision
artifact.

This paper makes a sharper claim, scoped to the $\ge 4$-bit formats that current
Blackwell-class hardware actually serves. \textbf{At these bit-widths the
expensive part of the cross-layer apparatus is unnecessary: a single additive
per-Linear cost, if it is faithful enough and measured at adequate numerical
precision, ranks allocations well enough that a lightweight real-KL selector
does the rest.} Our contributions:

\begin{enumerate}[leftmargin=*,itemsep=2pt]
  \item \textbf{\textsc{Aura}} (Section~\ref{sec:aura}): an $O(N)$
    production-faithful per-(Linear, format) cost, equal to a randomized
    KL/Gauss--Newton--Fisher estimate of the downstream KL contribution of the
    production-rendered weight error the exporter ships (stored bf16;
    Section~\ref{sec:limits}). It improves served KL by $17.9\%$ at 27B (matched
    bpp, two corpora) against the strong Fisher-weighted baseline that was the
    production default at measurement time, and by $38\text{--}40\%$ at the 4B
    knee against a weaker output-MSE baseline; Section~\ref{sec:experiments}.
  \item \textbf{A measurement-artifact study} (Section~\ref{sec:additivity})
    that explains why an additive cost suffices: a constant per-measurement
    bf16 floor inflates apparent non-additivity; in fp32 ($\equiv$ fp64) the
    unary KLs are nearly additive. We label this \emph{evidence}, not proof, and
    scope it to $\ge 4$ bits.
  \item \textbf{The operational rate--distortion geometry}
    (Section~\ref{sec:rd}): the achievable frontier is the integral of the
    sorted cost-per-bit, hence convex; the sharp knee we observe is induced by a
    discontinuity in the menu's marginal rates (an FP8 rung), not by dispersion in
    raw sensitivity.
    This corrects a natural misreading of an apparently ``linear'' frontier and
    reframes selection as byte-budget / saturation rather than knee detection.
\end{enumerate}

We are deliberate about what \textsc{Aura} is \emph{not}. Its estimator is a
clean application of known second-order machinery (HAWQ-style Fisher sensitivity,
randomized Hessian sketches, KL--Fisher equivalences). We do not claim new
sensitivity theory; we claim a production-faithful realization plus a
measurement correction, validated on the serving metric. And we are explicit
(Sections~\ref{sec:experiments}--\ref{sec:limits}) that the headline is a
KL-vs-BF16 win --- a distributional-fidelity metric, a proxy for reference
agreement and tool-call fidelity --- not a perplexity win, and that \textsc{Aura}'s weight-error cost
is not a complete model of W4A4 served error.

% =====================================================================
\section{Background: the \textsc{PrismaQuant} substrate}\label{sec:background}

\paragraph{Per-Linear format assignment.} We assume the \textsc{PrismaQuant}
serving substrate \cite{prismaquant}. A model is a set of Linears
$\Linears$; an assignment $\Assignment:\Linears\to\Formats$ gives each Linear one
format from a menu $\Formats\subseteq\{$BF16, FP8, NVFP4, $\dots\}$, subject to
serving constraints: fused siblings ($q/k/v$ in attention,
$\mathrm{gate}/\mathrm{up}/\mathrm{down}$ in MoE) must share a format, and
BF16/FP8\_SOURCE are passthrough-only (selectable only when the source tensor is
already that precision). The artifact is exported as stock
\texttt{compressed-tensors} and served by unmodified vLLM \cite{vllm} on
Blackwell-native kernels, with no forked runtime. Body bits-per-parameter
($\bpp$) is reported over \emph{quantizable} Linears only (excluding
\texttt{lm\_head}, MTP and visual sidecars).

\paragraph{Formats.} NVFP4 \cite{nvfp4_blackwell} is a 4-bit element format with
FP8 group scales over groups of 16 ($4 + 8/16 = 4.5$\,bpp weight-only), served
as W4A4 with group-16 FP4 activation quantization; FP8 (E4M3) is the 8-bit rung,
served W8A8 with dynamic per-token activation scales and \emph{no} new kernel;
BF16 is lossless passthrough. The microscaling MXFP family \cite{ocp_mx} is
available but de-menued for inference, since exact-scale FP8 Pareto-dominates
MXFP8 on already-served layers.

\paragraph{The shipping criterion.} A candidate ships only if it wins on the
\emph{serving metric}: exact full-vocab vLLM KL-vs-BF16 on a held-out split,
plus a perplexity / tail-NLL gate on the served artifact. The held-out split is
disjoint from the text used to generate any cost. This criterion is the contract
the rest of the paper is measured against.

\paragraph{Allocation as a knapsack.} Given a per-(Linear, format) cost
$c_{i,f}\ge 0$ and a bit count $b_{i,f}$ (exact rendered bytes including scale
overhead) with parameter count $n_i$, the additive allocator solves the
multiple-choice knapsack
\begin{equation}\label{eq:mckp}
  \min_{\Assignment}\ \sum_{i\in\Linears} c_{i,\Assignment_i}
  \quad\text{s.t.}\quad
  \sum_{i\in\Linears} b_{i,\Assignment_i}\, n_i \;\le\; B\!\sum_{i\in\Linears} n_i,
\end{equation}
by an exact bin-indexed dynamic program, with fused siblings grouped into joint
choices. Equation~\eqref{eq:mckp} is the entire optimizer; the paper is about
what $c_{i,f}$ should be, and about the fact that its solutions are
\emph{generated}, then \emph{selected} by real KL.

\paragraph{Related work.} Beyond CLADO/HAWQ-V3/CoopQ \cite{clado2023,hawqv3,coopq}
and the AutoML allocator AMQ \cite{amq}, the cost \textsc{Aura} uses belongs to a
well-developed family. Orr, Ribar and Luschi \cite{orr2025} derive Fisher--KL
additive per-tensor format allocation --- essentially the $\frac12 H\!\cdot\!
\mathrm{MSE}$ surrogate that is \textsc{PrismaQuant}'s own $L_1$ baseline.
FIMA-Q \cite{fimaq2025} builds a KL-proportional Fisher-information quant loss
for vision transformers; ``A KL Lens on Quantization'' \cite{kllens2026} uses
per-layer KL sensitivity for mixed precision but forward-only (backprop-free).
\textsc{Aura}'s narrow delta over this family is (i) a randomized
weight-gradient projection of the downstream KL--Fisher quadratic, (ii) measured
on the \emph{production-rendered} (not RTN) error, (iii) over a multiple-choice
\emph{format} knapsack, (iv) selected on served KL. We make the novelty scoping
explicit in Section~\ref{sec:limits}.

% =====================================================================
\section{The measurement-artifact study}\label{sec:additivity}

The additive objective in \eqref{eq:mckp} is, for any per-Linear cost that holds
the other Linears fixed, a first-order expansion of $\kl(\Assignment)$. The
cross-layer literature --- and our own prior cascade --- is motivated by the
observation that this objective \emph{overshoots}: committing many flips at once
steps outside the trust region, and predicted decreases become measured
increases \cite{clado2023,prismaquant}. We report evidence that, at $\ge 4$ bits,
much of this overshoot is not cross-layer coupling at all but a numerical-%
precision artifact of how the unary KLs are measured.

\paragraph{A constant per-measurement floor reproduces the ``non-additivity''
signature.} Suppose each measured unary KL carries, in addition to its true
value, a constant additive floor $\varepsilon$ independent of the Linear (the
residual of differencing two near-equal categorical distributions at finite
precision), and that the true per-Linear KLs are additive.

\begin{proposition}[Floor-induced pseudo-coupling]\label{prop:floor}
Let $\widehat{\kl}(S) = \kl^\star(S) + \varepsilon$ for every measured subset
$S\subseteq\Linears$, where $\kl^\star$ is exactly additive,
$\kl^\star(S)=\sum_{i\in S}\kappa_i$, and $\varepsilon>0$ is a constant. Then for
a set of $m$ singletons,
\[
  \frac{\sum_{i=1}^m \widehat{\kl}(\{i\})}{\widehat{\kl}(\{1,\dots,m\})}
  \;=\; \frac{\sum_i \kappa_i + m\varepsilon}{\sum_i\kappa_i + \varepsilon},
\]
which exceeds $1$ by $(m-1)\varepsilon/(\sum_i\kappa_i+\varepsilon)$, and every
pairwise interaction is the \emph{same} negative constant,
\[
  I(i,j) \;\coloneqq\; \widehat{\kl}(\{i,j\}) - \widehat{\kl}(\{i\}) - \widehat{\kl}(\{j\})
  \;=\; -\varepsilon,
\]
independent of $\kappa_i,\kappa_j$.
\end{proposition}

\noindent\emph{Proof.} Direct substitution; the $\varepsilon$ terms count
$m$ in the numerator and $1$ in the denominator, and
$I(i,j)=(\kappa_i+\kappa_j+\varepsilon)-(\kappa_i+\varepsilon)-(\kappa_j+\varepsilon)
=-\varepsilon$. \hfill$\square$

\medskip
Proposition~\ref{prop:floor} is falsifiable: it predicts (a) a sum/joint ratio
that inflates with the number of Linears even with zero true coupling, and (b)
\emph{all-negative, magnitude-independent} pairwise interactions. Both
signatures appear, and vanish with precision.

\paragraph{Precision sweep (Qwen3-4B, weight-only W4A16, full-vocab
KL computed in-dtype, 12 Linears / 6 pairs).} Table~\ref{tab:precision} reports
the sum/joint ratio and mean interaction as the KL arithmetic is moved from bf16
to fp32 to fp64. In bf16 the ratio is $1.88$ with all-negative interactions ---
the floor signature. In fp32 it collapses to $0.92$ with interactions two orders
of magnitude smaller, and \textbf{fp64 is identical to fp32 to five significant
figures}, so the floor is a bf16 artifact and higher precision than fp32 is
moot. A small \emph{genuine} residual remains (the $0.92$ is mildly
super-additive --- the joint slightly exceeds the unary sum, i.e.\ the additive
surrogate slightly \emph{under}-counts), is precision-independent, and is roughly
$10\times$ smaller than the bf16 floor.

\begin{table}[t]\centering
\small
\caption{The apparent non-additivity is a bf16 measurement floor. Qwen3-4B,
W4A16, same 12 Linears / 6 pairs, KL computed in the stated dtype. $\Sigma/J$ is
sum-of-unary over joint; $\bar I$ is the mean pairwise interaction.
Repro: \texttt{xlayer\_rung0/fp64\_additivity.py} (run dir since archived; values
were re-derived verbatim from its raw logs in the 2026-06-12 external review).}
\label{tab:precision}
\begin{tabular}{lccl}
\toprule
KL arithmetic & $\Sigma/J$ & $\bar I$ & interaction sign \\
\midrule
bf16 & $1.884$ & $-5.27\times10^{-4}$ & all negative (floor) \\
fp32 & $0.919$ & $-1.22\times10^{-5}$ & scattered $\sim 0$ \\
fp64 & $0.919$ & $-1.22\times10^{-5}$ & $\equiv$ fp32 (5 s.f.) \\
\bottomrule
\end{tabular}
\end{table}

\paragraph{The signature reproduces at scale (Qwen3.6-27B, served-%
faithful W4A4, bf16 forward / fp32 KL).} On 12 Linears of the 27B model,
$\sum$(unary)$=8.53\times10^{-3}$ against a true joint of $3.86\times10^{-3}$, a
$2.21\times$ overshoot; all 12 \emph{sampled} pairwise interactions are negative with mean
$-3.7\times10^{-4}$ and standard deviation $1.3\times10^{-4}$ across a $9\times$
span of unary magnitudes. A single constant $\varepsilon\approx4.2\times10^{-4}$
explains \emph{both} the overshoot ($\approx 11\varepsilon$, matching
Proposition~\ref{prop:floor} with $m=12$) and the constant interaction
($\approx-\varepsilon$). Constancy across a $9\times$ magnitude span is the
discriminator: a genuine second-order coupling would scale with the
perturbations, a measurement floor does not.

\begin{remark}[Scope, stated honestly]\label{rem:scope}
This is evidence, not proof, and it is scoped to $\ge 4$ bits. The clean
fp32-vs-fp64 contrast is at 4B (a resident fp32 27B forward is $108$~GB and does
not fit our $128$~GB box); at 27B we show the \emph{mechanism} (constant-%
$\varepsilon$ fit) on served-faithful W4A4. At $2$--$3$ bits the genuine coupling
grows --- perturbations $\sim\!4\times$ larger make second-order terms
$\sim\!16\times$ larger --- and CoopQ \cite{coopq} reports its largest gains
exactly there ($2.0$--$2.5$~bits). We therefore do \emph{not} claim
``non-additivity is a measurement artifact'' in general; we claim that in the
$\ge 4$-bit serving regime an additive cost, measured in fp32, plus a real-KL
selector, is empirically sufficient, and that the bf16 floor explains why the
overshoot looked larger than the genuine residual it must absorb.
\end{remark}

\paragraph{What this retires.} The prior \textsc{PrismaSCOUT} system built an
$L_2$ perturbed-activation fixed point and an $L_3$ propagated end-KL pass
specifically to \emph{model} the cross-layer interaction that
Proposition~\ref{prop:floor} and Table~\ref{tab:precision} now attribute, in the
$\ge 4$-bit regime, mostly to a bf16 measurement floor. We therefore retire the
$L_2/L_3$ cost modeling. We do \emph{not} retire the principle it served ---
``surrogates generate, real measurement selects'' --- which is exactly what
absorbs the small genuine residual: the real-KL selector is kept, and the
\emph{cost} side collapses to one fp32 additive pass (Section~\ref{sec:aura}).
We are not claiming the cascade was without value --- it produced the shipped
flagship artifact --- only that its motivating quantity is, at these bit-widths,
largely an artifact of how it was measured, so a faithful per-Linear cost plus
selection is the simpler and equally-served path forward. A matched-bpp
head-to-head at 4B makes this concrete (Section~\ref{sec:experiments}): the full
$L_2$ perturbed-X cascade improves on the $L_1$ baseline by only $1.5\%$
confident-KL, while \textsc{Aura} improves by $39\%$ --- so \textsc{Aura} beats
the cascade by $38.5\%$, and the cross-layer modeling the cascade performs buys
almost nothing. ($L_3$ polish was archived and not run; this is the $L_2$ fixed
point, the cascade's core cross-layer cost.)

% =====================================================================
\section{\textsc{Aura}: a production-faithful KL--Fisher cost}\label{sec:aura}

\paragraph{The downstream-KL contribution of one Linear.} Quantizing Linear $i$
to format $f$ perturbs its weight by $\dW_{i,f}=Q_f(W_i)-W_i$, displacing the
model's logits and hence the output distribution. To second order, the increase
in forward KL from the BF16 reference is the Fisher/Gauss--Newton quadratic of
the induced logit displacement. Writing $J_i$ for the Jacobian of the logits
with respect to $W_i$ and $F$ for the categorical Fisher of the reference
distribution, the unary contribution is
$\tfrac12\,\mathrm{vec}(\dW_{i,f})^\top (J_i^\top F J_i)\,\mathrm{vec}(\dW_{i,f})$.

\begin{definition}[\textsc{Aura} cost]\label{def:aura}
Let $g^{(k)}_{W_i}=\nabla_{W_i}\,\phi_k(\text{logits})$ be the weight gradient of
the $k$-th randomized KL--Fisher probe $\phi_k$, a Rademacher logit functional
\emph{constructed} so that its covariance equals the categorical Fisher
normalized over the selected token positions (a property of the probe
construction, not of the gradient; cf.\ \cite{kllens2026}), backpropagated through
the \emph{clean} (full-precision) model. The \textsc{Aura} cost is
\begin{equation}\label{eq:aura}
  c_{i,f} \;=\; \tfrac12\,\frac{1}{K}\sum_{k=1}^{K}
     \big\langle\, g^{(k)}_{W_i},\; \dW_{i,f}\,\big\rangle^2 .
\end{equation}
\end{definition}

\begin{proposition}[Unbiasedness]\label{prop:unbiased}
If $\E_k\!\big[\,\mathrm{vec}(g^{(k)}_{W_i})\,\mathrm{vec}(g^{(k)}_{W_i})^\top\big]
= J_i^\top F J_i$, then $\E_k[c_{i,f}] = \tfrac12\,
\mathrm{vec}(\dW_{i,f})^\top(J_i^\top F J_i)\,\mathrm{vec}(\dW_{i,f})$, the exact
second-order unary KL contribution.
\end{proposition}

\noindent\emph{Proof.} $\E_k\langle g^{(k)},\dW\rangle^2
=\mathrm{vec}(\dW)^\top\E_k[\mathrm{vec}(g^{(k)})\mathrm{vec}(g^{(k)})^\top]
\mathrm{vec}(\dW)$. \hfill$\square$

\medskip
Proposition~\ref{prop:unbiased} is conditional on the covariance identity above,
which is a property of how $\phi_k$ is constructed (the standard randomized-Fisher
probe), not something we prove here; what we validate end-to-end is the served-KL
win (Section~\ref{sec:experiments}), not the estimator's exactness.
Equation~\eqref{eq:aura} is $O(N)$: $K$ backward passes through the clean model
yield every $g^{(k)}_{W_i}$ at once, and each cost is a single inner product
against the format's rendered error. There is no quantized forward pass and no
pairwise term --- the total cost summed in \eqref{eq:mckp} is the additive
surrogate whose soundness Section~\ref{sec:additivity} argues for.

\paragraph{Three things that make it production-faithful.}
(1)~\emph{Production-rendered error.} $\dW_{i,f}$ is read from the
\textsc{ProductionWeightCache} --- the same GPTQ\,+\,JSO\,+\,static-act-order
rendering the exporter ships (stored bf16, and falling back to RTN on a cache
miss; Section~\ref{sec:limits}), not a fresh RTN round. The cost therefore scores
the rendered artifact, not an idealization of it. (2)~\emph{Format table.}
$c_{i,f}$ is computed for every lossy $f$ in the menu (NVFP4, FP8), with
BF16/FP8\_SOURCE hardwired to zero (passthrough), so \eqref{eq:mckp} chooses
among real serving rungs. (3)~\emph{Selection on served KL.} The allocator's
Pareto candidates are re-scored by exact vLLM KL-vs-BF16 on the held-out split;
\textsc{Aura} is the \emph{generator}, real KL is the \emph{selector}, exactly as
Section~\ref{sec:additivity}'s residual coupling requires.

\paragraph{Scale without surgery.} \textsc{Aura} is a signal-dominated quadratic
$\tfrac12\langle g,\dW\rangle^2$, not a KL \emph{difference}, so it has no
cancellation floor and is safe to measure on a bf16-resident model: the
bf16-resident cost reproduces the fp32 cost with Spearman $0.9967$ and median
ratio $0.990$. A resident bf16 forward/backward fits a 27B model ($\sim$54~GB)
on the $128$~GB box, so no streaming probe surgery is needed
(Section~\ref{sec:limits} notes the one place bf16 \emph{does} enter the cost).

% =====================================================================
\section{Operational rate--distortion geometry}\label{sec:rd}

A natural question, once allocation is an additive knapsack, is the shape of the
loss-vs-bitrate frontier. The operational rate--distortion view \cite{ortega1998}
makes the answer exact, and corrects a tempting misreading.

\paragraph{Setup.} With a two-rung menu $\{$NVFP4, BF16$\}$, each Linear either
stays BF16 (lossless, full bits) or drops to NVFP4, saving
$\Delta b_i$ bits at additive cost $c_i$. Define the \emph{cost-per-bit}
$\rho_i = c_i/\Delta b_i$.

\begin{proposition}[Convex frontier from sorted cost-per-bit]
\label{prop:rd}
Consider the \emph{continuous} (fractional) relaxation of the additive knapsack.
Its minimum total cost at each bit budget is achieved by dropping Linears to
NVFP4 in increasing order of $\rho_i$, giving
$D(R)=\int_0^{B(R)}\rho_{\mathrm{sorted}}(b)\,db$, which is convex (its marginal
slope is the nondecreasing $\rho_{\mathrm{sorted}}$) and affine iff $\rho_i$ is
constant. The integer-optimal frontier lies on or above this convex envelope, and
the empirical frontier we plot is the non-dominated lower staircase of measured
$(\bpp,\kl)$ points.
\end{proposition}

\noindent\emph{Proof.} Let $B=B(R)$ be the saved-bit mass the budget demands and
write the relaxation over fractions $x_i\in[0,1]$: minimize $\sum_i x_ic_i$
subject to $\sum_i x_i\Delta b_i\ge B$. Order the Linears
$\sigma(1),\dots,\sigma(N)$ so that $\rho_{\sigma(1)}\le\dots\le\rho_{\sigma(N)}$,
and let $\rho_{\mathrm{sorted}}$ be the associated step function,
$\rho_{\sigma(k)}$ on $[\sum_{j<k}\Delta b_{\sigma(j)},\sum_{j\le k}\Delta b_{\sigma(j)})$.
(i)~\emph{The greedy fill is optimal.} Let $x$ be optimal and suppose there were
$p,q$ with $\rho_p<\rho_q$, $x_q>0$, $x_p<1$. For $\delta>0$ small replace
$x_q\leftarrow x_q-\delta/\Delta b_q$ and $x_p\leftarrow x_p+\delta/\Delta b_p$:
saved bits are conserved, so feasibility is kept, while cost changes by
$\delta(\rho_p-\rho_q)<0$, contradicting optimality. Hence every optimum takes an
ascending-$\rho$ prefix with any remainder confined to a single $\rho$-class
(interchanges within a class leave the cost $\rho\cdot b$ unchanged), of value
$D(B)=\int_0^B\rho_{\mathrm{sorted}}(b)\,db$.
(ii)~\emph{Convexity.} For $u<v<w$ the chord slopes are running means of a
nondecreasing integrand,
$\tfrac{D(v)-D(u)}{v-u}=\operatorname{avg}_{[u,v]}\rho_{\mathrm{sorted}}\le
\operatorname{avg}_{[v,w]}\rho_{\mathrm{sorted}}=\tfrac{D(w)-D(v)}{w-v}$,
the three-chord condition for convexity on $[0,\sum_i\Delta b_i]$; composing with
the affine map $R\mapsto B(R)$ preserves it.
(iii)~\emph{Affine iff constant.} If all $\rho_i=\rho$ then $D(B)=\rho B$. Conversely
$D$ is absolutely continuous with $D'(b)=\rho_{\mathrm{sorted}}(b)$ a.e.; if $D$ is
affine then $D'$ is a.e.\ a constant $\bar\rho$, and a step function that is a.e.\
constant is constant on each of its steps (every step has positive width), so every
$\rho_i=\bar\rho$.
(iv)~\emph{Integer versus fractional.} Any assignment is a feasible $x$, so the
integer optimum lies weakly above $D$; with unequal $\Delta b_i$ the gap can be
strict, since the fractional solution splits one item where integrality must
overpay. The plotted object is neither of these but the non-dominated lower
staircase of measured $(\bpp,\kl)$ points: sampled at finitely many budgets, it sits
weakly above the integer frontier it samples and is convex only through its lower
convex hull. \hfill$\square$

\begin{proposition}[The knee tracks the cost-per-bit density, not raw sensitivity]\label{prop:knee}
Work in the fractional relaxation of Proposition~\ref{prop:rd} with several lossy
formats: each Linear contributes the segments of the lower convex hull of its rung
points $(0,0)$ and $(\beta_{i,f},c_{i,f})$, where $\beta_{i,f}$ is the bits BF16
loses to format $f$; a segment saves $\beta$ bits at cost $c$, ratio
$\rho=c/\beta$, and the greedy fill runs over all segments jointly. Then
(i)~$D$ is piecewise affine --- affine along each maximal run of equal sorted
ratio --- and every slope break, hence every knee, occurs at a gap between distinct
sorted ratios, of size equal to that gap; dispersing the raw costs $\{c_i\}$ while
holding ratios fixed (e.g.\ $c_i=\rho\,\Delta b_i$ with varied $\Delta b_i$) leaves
the frontier exactly affine. (ii)~A separated rung is sufficient for a sharp knee:
if some rung $f$ satisfies $\max_i\rho_{i,f}<\min_i\rho_{i,g}$ for another rung
$g$, the entire $f$-block is consumed before any $g$-segment, and the frontier's
slope drops by at least $\min_i\rho_{i,g}-\max_i\rho_{i,f}>0$ at the byte count
where the cheaper rung is exhausted (for FP8: BF16$\to$FP8 buys many bits at
near-zero cost, FP8$\to$NVFP4 is expensive).
\end{proposition}

\noindent\emph{Proof.} (i) Proposition~\ref{prop:rd}'s exchange argument applies
verbatim to the segment list, so $D(B)=\int_0^B\rho_{\mathrm{sorted}}(b)\,db$ over
sorted segment ratios; $D$ is affine exactly on the fill intervals of maximal
constant-ratio runs, and crossing a boundary between ratios $\rho'<\rho''$ changes
the slope by $\rho''-\rho'$. The raw-cost clause is (iii) of
Proposition~\ref{prop:rd} read backwards. (ii) Separation places the whole
$f$-block first in sorted order; crossing the block boundary drops the slope by
$\min_i\rho_{i,g}-\max_i\rho_{i,f}>0$, a single jump of that size. \hfill$\square$

\medskip
Two empirical clauses sit deliberately outside the proposition. Under a single
lossy format no separated rung exists; whether the resulting curve bends smoothly
is a property of the joint distribution of $(c_i,n_i)$ --- on our models
$\rho_i\propto c_i/n_i$ varies smoothly and the two-rung frontier is visibly convex
with no interior slope break. And a sufficiently heavy-tailed $\{\rho_i\}$ would
produce a knee with no help from the menu, since (i) is indifferent to the origin
of the gap; we observe no such gap on the two-rung menu here.

\paragraph{Consequences, on real data.} On the Qwen3.6-27B
\textsc{Aura} costs, the cost-per-bit is far from democratic:\footnote{Derivation
status of the summary statistics: the Gini coefficients, the $\rho$ ratio, the
chord curvatures quoted in this section, and the bf16-vs-fp32 rank/median
agreements of Section~\ref{sec:aura} were computed once during development and
carried forward in working notes; unlike the served tables of
Section~\ref{sec:experiments}, they have not been re-derived from retained raw
logs, and their scratch run directories are archived. Read them as indicative
magnitudes, not audited results.} weighted Gini
$0.45$ (unweighted $0.65$), with $\rho_{\max}/\rho_{\mathrm{median}}\approx 63$.
By Proposition~\ref{prop:rd} the frontier is therefore convex, and the
reconstructed envelope (integral of sorted $\rho$) has $33\%$ \emph{chord
curvature} --- the maximum vertical deviation of the curve from its straight chord
as a fraction of the KL range --- over $4$--$16$~bpp. The measured frontier --- resident W4A4 emulation, a served-faithful
screen, \emph{not} a served-vLLM measurement --- is convex in the same sense
and knees in the same region ($24.5\%$ chord curvature, knee
$\approx6.5$~bpp); the smaller magnitude is expected for a single-seed
screen against a $496$-Linear surrogate reconstruction, and we do not read
the two as quantitatively matched. It \emph{looks} linear only because production plots
$\log_{10}\kl$ over the narrow $4.5$--$6.0$~bpp window where allocations live ---
an exponential $\kl\!\sim\!10^{-aR}$ is straight in log-KL by construction
(Figure~\ref{fig:aura_rd_geometry}b). By Proposition~\ref{prop:knee} a sharp
knee appears only when the FP8 rung is added (Figure~\ref{fig:aura_rd_geometry}a):
on Qwen3-4B the two-rung frontier is near-linear while
$\{$NVFP4, FP8, BF16$\}$ bends to a knee near $8$~bpp ($6\times$ lower KL at
matched bytes in a resident RTN weight-only train-split screen, not a served
held-out result). Crucially, a better cost model (\textsc{Aura} vs.\ the baseline)
shifts the frontier \emph{down} at fixed bpp --- a vertical translation
orthogonal to its curvature --- so ``\textsc{Aura} is the better ranker'' and
``the two-rung frontier looks linear'' are simultaneously true. A wider validation
corroborates the picture beyond the production window: on Qwen3.8-27B dense,
eleven held-out points spanning $4.50$--$8.25$~bpp fall on a single log-linear law,
$\log_{10}\kl=-0.1133-0.2357\,\bpp$ ($R^2=0.9948$), with no saturation in range
(Repro: \texttt{docs/design/operating\_point\_selection.md}) --- the narrow-window
linearity of Figure~\ref{fig:aura_rd_geometry}b is the visible part of a globally
exponential frontier. An exponential buys the same \emph{relative} return at every
bpp ($10^{0.2357}\!\approx\!1.72\times$ KL per bpp), so no interior operating point
is distinguished and a knee detector must manufacture one: on that same sweep
kneedle returns $4.75$~bpp spelled on the log axis and $6.00$~bpp spelled on the
linear axis. The $\approx6.5$-bpp knee above therefore survives as chord geometry
of a convex curve, not as an anchor. We therefore
select ship points by exact byte budget / raw-KL saturation, not by an axis-%
dependent knee detector \cite{kneedle2011}.

\input{figures/fig_aura_rd_geometry.tex}

% =====================================================================
\section{Experiments}\label{sec:experiments}

All comparisons hold the rendering fixed --- one shared GPTQ+JSO production cache
for both arms, so only the \emph{allocation} differs --- and report exact vLLM
KL-vs-BF16 on a held-out split disjoint from cost generation (the
selection/validation split; we return to the absence of a third, untouched test
split in Section~\ref{sec:limits}). The two scales use \emph{different} baselines,
which we state plainly: the 27B comparison is against the Fisher-trace$\times$%
output-MSE cost
($u_i=\tfrac12\,\mathrm{Tr}(\hat H_i)\,\mathrm{MSE}^{\mathrm w}_i$), a strong
Fisher-weighted local surrogate and the production default at the time of these
measurements;\footnote{The production default has since moved to
$\mathrm{Tr}(\hat H_i)\,\mathrm{MSE}_{W_i}$ (weight-MSE): a subsequent served A/B
found the output-MSE term double-counts activation energy already present in the
Fisher trace, decisively at small scale ($-51\%$/$-58\%$ pooled KL at 4B/0.6B over
five calibration draws) and neutral-to-positive at 27B. An \textsc{Aura}-vs-%
weight-MSE-baseline head-to-head has not been run at 27B; the $17.9\%$ here is
against the output-MSE variant, stated exactly.} the 4B comparison is against the
weaker raw output-MSE cost. The larger 4B gap is therefore partly the weaker
baseline --- the airtight, fair-baseline result is the 27B one.

\paragraph{Qwen3-4B, at the knee.} \textsc{Aura}'s sensitivity is
\emph{concentrated} (a few MLP gate/up Linears carry the KL; sharp kneedle at
$\approx4.84$~bpp); the baseline is \emph{diffuse} (protects many late-attention
Linears, fuzzy knee $5.2$--$6.0$). At the bit-starved knee the difference is
large (Table~\ref{tab:4b}); above it both have banked the load-bearing
protections and the gap washes out --- the operative lesson is to test at the
knee, not above it.

\begin{table}[t]\centering
\small
\caption{Qwen3-4B, served vLLM KL-vs-BF16 at the $4.85$~bpp knee (lower is
better), shared NVFP4+BF16 cache, two independent calibrations. Baseline here is
the \emph{weaker} raw output-MSE cost (cf.\ the strong Fisher-weighted baseline
used at 27B), so the gap is not directly comparable to Table~\ref{tab:27b}.
WikiText PPL at the same point. Repro: \texttt{xlayer\_rung0/} (since archived;
values verified against its raw artifacts in the 2026-06-12 review).}
\label{tab:4b}
\begin{tabular}{lccc}
\toprule
 & \textsc{Aura} & output-MSE baseline & $\Delta$ \\
\midrule
confident KL, rep1 ($n{=}4$, $K{=}16$) & $0.2268$ & $0.3645$ & $-37.8\%$ \\
confident KL, rep2 ($n{=}8$, $K{=}32$) & $0.2204$ & $0.3645$ & $-39.5\%$ \\
WikiText PPL (BF16 ref.\ $3.582$)      & $3.965$  & $4.834$  & $-18\%$ \\
\bottomrule
\end{tabular}
\end{table}

\paragraph{Qwen3.6-27B, matched $5.99$~bpp, two corpora.} The win holds at scale
against the \emph{strong} baseline. The gap is larger at $5.99$ than at $5.25$~bpp
($-17.1\%$ vs.\ $-4.2\%$ confident KL on the small corpus): the baseline plateaus
($0.0318\!\to\!0.0319$ from $5.25\!\to\!6.0$~bpp; $+39$ BF16 Linears bought
nothing) while \textsc{Aura} keeps improving ($0.0305\!\to\!0.0265$). Both
trajectories are single-seed and lack confidence intervals
(Section~\ref{sec:limits}); we report the trend, not a significance claim. The
relative advantage \emph{is} stable across two disjoint held-out corpora of
different size and content ($-17.1\%$ and $-17.9\%$) --- this cross-corpus
replication is our main evidence that it is not a single-split screen.
\textsc{Aura} spends its budget concentratedly ($341$ NVFP4 / $155$ BF16, on the
large MLPs) where the baseline is diffuse ($314$ / $182$).

\paragraph{A held-out test split.} The costs are generated on the WikiText
\emph{train} split, so to rule out a selection screen we re-served both arms'
exported artifacts and measured served vLLM confident-KL on the WikiText
\emph{test} split --- text disjoint from cost generation. The advantage holds
and, if anything, strengthens: $\mathbf{-39.4\%}$ confident-KL at 4B
(\textsc{Aura} $0.245$ vs.\ $0.405$; $-37.7\%$ all-position) and
$\mathbf{-33.5\%}$ at 27B (\textsc{Aura} $0.0344$ vs.\ $0.0517$; $-22.1\%$
all-position). The 27B confident-KL gap is therefore $-17$ to $-34\%$ across
two disjoint held-out corpora, with the $-33.5\%$ point a second WikiText-test
measurement window rather than an independent third corpus --- always
substantial, always favouring \textsc{Aura}. (Caveat for our own tooling: a
resident HF W4A4 \emph{emulation}
of the same allocations badly under-counts the baseline's penalty at 4B versus
the vLLM CUTLASS kernel; the served numbers here are the authority.)

\paragraph{Head-to-head against the cascade.} The sharpest test of the retirement
(Section~\ref{sec:additivity}) puts \textsc{Aura} directly against the full $L_2$
perturbed-X cascade it replaces, at matched bpp on the same production cache,
exported and served identically (Table~\ref{tab:cascade}). The cascade --- the
expensive cross-layer cost the prior system paid for --- improves on the $L_1$
Fisher-weighted baseline by only $1.5\%$ confident-KL; \textsc{Aura}'s additive
KL-Fisher cost improves by $39.4\%$, beating the cascade by $38.5\%$.
Mechanistically the baseline and cascade both diffusely protect late-attention
$v$-projections, while \textsc{Aura} concentrates bits on the KL-load-bearing MLP
$\mathrm{gate}/\mathrm{up}$; the $L_2$ fixed point refines \emph{within} the
baseline's strategy rather than discovering \textsc{Aura}'s. ($L_3$ polish modules
were archived and not run; this is the $L_2$ fixed point --- the cascade's core
cross-layer cost --- re-allocated through the same allocator and fused-sibling
promotion as the other arms, so only the per-Linear cost differs.)

\begin{table}[t]\centering
\small
\caption{Qwen3-4B, served vLLM confident-KL-vs-BF16 on the held-out WikiText test
split, matched $\sim\!4.83$~bpp, one production cache and one export/serve path for
all three arms (only the per-Linear cost differs). The full $L_2$ cascade barely
beats $L_1$; \textsc{Aura} beats both by $\sim\!38\%$. Repro: \texttt{aura-blocker-runs/}.}
\label{tab:cascade}
\begin{tabular}{lccc}
\toprule
per-Linear cost & bpp & confident-KL & all-position KL \\
\midrule
$L_1$ Fisher$\times$output-MSE (baseline) & $4.85$ & $0.405$ & $0.463$ \\
$L_2$ perturbed-X fixed point (cascade)   & $4.83$ & $0.399$ & $0.459$ \\
\textsc{Aura} (KL-Fisher adjoint)         & $4.84$ & $\mathbf{0.245}$ & $\mathbf{0.289}$ \\
\bottomrule
\end{tabular}
\end{table}

\begin{table}[t]\centering
\small
\caption{Qwen3.6-27B, served vLLM KL-vs-BF16 at matched $5.99$~bpp (NVFP4+BF16),
one shared cache both arms. ``conf.'' = confident-position mean KL.
Repro: \texttt{aura-27b-flagship/}.}
\label{tab:27b}
\begin{tabular}{lccc}
\toprule
held-out corpus & \textsc{Aura} & Fisher$\times$MSE & $\Delta$ (conf.) \\
\midrule
small (798 conf.\ pos.)            & $0.02646$ & $0.03193$ & $-17.1\%$ \\
WikiText-test (4393 conf.\ pos.)   & $0.05594$ & $0.06815$ & $-17.9\%$ \\
\midrule
\multicolumn{4}{l}{\emph{all-position} KL (WikiText-test): $0.06969$ vs $0.07899$, $-11.8\%$; top-1 $+0.57$pp}\\
\bottomrule
\end{tabular}
\end{table}

\paragraph{The honest perplexity nuance.} The KL win does \emph{not} carry to raw
perplexity at 27B. At $5.99$~bpp (WikiText-test, 8176 tokens): BF16 $9.243$,
\textsc{Aura} $9.460$ ($+2.35\%$), baseline $9.496$ ($+2.74\%$) --- \textsc{Aura}
is better by only $0.38\%$ relative ($0.39$ percentage points less degradation).
At $6$~bpp perplexity is near-saturated for both arms, while KL is not (Aura still
improving, baseline plateaued). The 27B result is therefore specifically a
\textbf{distributional-fidelity} win --- a lower KL-vs-BF16, a proxy for reference
agreement and tool-call fidelity --- not a
perplexity win. Perplexity does not veto (Aura $\le$ baseline), but we do
not claim it as a PPL improvement. The $-18\%$ PPL gap at 4B was against a
\emph{weaker} (no-Fisher) baseline and should not be read as the 27B story.

\paragraph{Behavioral corroboration (tool-call fidelity).} The
distributional-fidelity framing now has a direct behavioral check. On a
deterministic tool-use benchmark (ToolEvalBench, hard mode, sequential,
temperature $0$, one fixed seed, same harness for every row), an
\textsc{Aura}-allocated 27B artifact at matched $5.5$~bpp scores $88/100$
against $84/100$ for the prior production allocation at the same bpp; the
shipped \textsc{Aura} artifact (same $5.5$~bpp point, corrected render/export
code, full NVFP4/FP8/BF16 menu) scores $91/100$ against $86/100$ for the BF16
reference itself and $85/100$ for the prior $5.31$-bpp flagship. Two caveats,
stated plainly: these are \emph{artifact-level} comparisons (menu and code
version vary alongside the cost, unlike the matched-cache A/Bs above), and a
quantized model scoring above BF16 on a single deterministic run should be read
as \emph{parity within run-to-run structure}, not superiority. What the result
does support is the claim the KL metric is a proxy for: the \textsc{Aura}
allocations preserve tool-call behavior at least as well as the allocations they
replaced.

\paragraph{Weight-only cost vs.\ W4A4 served error (attribution).} \textsc{Aura}
scores weight error only, yet NVFP4 ships as W4A4. On the 27B artifact a
weights-only resident screen (W4A16) recovers $0.0252$, only $45\%$ of the served
KL ($0.0559$); adding group-16 FP4 activation quantization (W4A4) recovers
$0.0455$, $81\%$ of served --- so activation quantization is the \emph{majority}
of NVFP4's error. The relevant fact for allocation is that the \emph{relative}
Aura-vs-baseline gap is preserved through the rendering: $-16.1\%$ at W4A16
$\approx$ $-17.9\%$ served. We therefore claim that \textsc{Aura}'s weight-error
ranking \emph{survives} W4A4 validation, not that it is a complete W4A4 error
model (Section~\ref{sec:limits}).

% =====================================================================
\section{Limitations and honest accounting}\label{sec:limits}

\begin{enumerate}[leftmargin=*,itemsep=3pt]
  \item \textbf{Weight-only cost under W4A4.} \textsc{Aura} models the weight
    error $\dW_{i,f}$; activation quantization is $\sim$half of NVFP4's served
    error (above). \textsc{Aura} is a production-weight allocation \emph{ranker}
    whose selected assignments survive served W4A4 validation, not a complete
    model of W4A4 distortion. An activation-aware cost term is future work.
  \item \textbf{The rendered error is stored bf16.} In the current
    implementation $\dW_{i,f}$ is materialized in bf16 before the inner product
    \eqref{eq:aura}; the probe gradients are fp32 but the error vector is
    rounded. ``fp32 cost'' thus refers to the KL/Fisher arithmetic, not to
    $\dW$. We report this rather than overstate it; the fp32-$\dW$ version (and a
    coverage report for production-cache hits vs.\ RTN fallback) is the first
    hardening item below.
  \item \textbf{fp64 purity is 4B-only.} The decisive bf16-vs-fp32-vs-fp64
    contrast (Table~\ref{tab:precision}) fits only at 4B; at 27B we show the
    mechanism, not the precision-pure ladder. A resident fp32 27B forward needs
    $108$~GB.
  \item \textbf{Single-seed corpora at 27B.} The 27B KL is single-seed per
    corpus; the cross-corpus replication is the stability evidence, but
    multi-seed CIs are not yet in hand. (The 4B win has tight error bars across
    two independent calibrations.)
  \item \textbf{Split protocol.} Costs are generated on the WikiText \emph{train}
    split. The headline \textsc{Aura}-vs-baseline numbers are matched-bpp
    \emph{cost-model} comparisons (neither arm's allocation is tuned on the eval
    split), and we additionally report served confident-KL on the held-out
    WikiText \emph{test} split (Section~\ref{sec:experiments}), where the win holds
    ($-39\%$ at 4B, $-33\%$ at 27B). What remains is a fully pre-registered
    three-way cost/selection/test protocol with the count of screened candidates
    stated.
  \item \textbf{Cascade head-to-head is $L_2$-only and 4B-only.} We \emph{do}
    report a matched-bpp head-to-head (Table~\ref{tab:cascade}, \textsc{Aura} beats
    the $L_2$ cascade by $38.5\%$), but on Qwen3-4B and with the $L_2$ fixed point
    only --- the $L_3$ polish modules are archived and were not run. A fuller
    cascade ($L_3$, a larger calibration than the box's memory caps allowed) and a
    27B repeat would strengthen the retirement further; on the present evidence the
    $L_2$ cross-layer cost adds only $1.5\%$ over $L_1$.
  \item \textbf{Novelty scoping.} The \textsc{Aura} estimator is a clean
    synthesis of known machinery: Fisher--KL additive allocation
    \cite{orr2025}, KL-proportional Fisher quant loss \cite{fimaq2025},
    per-layer KL sensitivity \cite{kllens2026}, and randomized second-order
    sketches. We claim the \emph{production-faithful realization} (rendered-%
    error scoring, format knapsack, served-KL selection) and the
    \emph{measurement-artifact correction}, not new sensitivity theory.
  \item \textbf{MoE routed experts are outside the smooth cost's reach.} The
    experiments here are dense models, and measurement after this draft shows
    that boundary is real, not incidental: on a $192$-expert top-$8$ MoE
    (Qwen3.6-35B-A3B), quantizing a routed expert perturbs the \emph{router},
    changing which experts execute --- an error mode invisible to any
    fixed-route weight-error quadratic, \eqref{eq:aura} included. Under the
    faithful production render the cost's rank correlation with measured
    per-expert-unit KL drops (Spearman $0.45\!\to\!0.35$), and predicted
    NVFP4-vs-FP8 cost ratios of $2$--$49\times$ compress to a measured
    $1.1$--$1.5\times$. The production remedy follows this paper's own spine
    --- when the surrogate cannot see the cost, measure it: packed routed
    experts take \emph{empirical} per-unit end-to-end KL costs, merged into the
    same knapsack \eqref{eq:mckp} alongside \textsc{Aura} costs for the
    non-expert Linears, with real-KL frontier selection unchanged. The hybrid
    ships (served confident KL $0.029$--$0.031$ at $4.66$--$4.75$~bpp on the
    35B), but the paper's additive-cost claim should be read as scoped to
    Linears that do not gate control flow.
\end{enumerate}

\paragraph{The experiments that would most harden this paper, in priority order.}
(1)~\emph{partly done} --- served confident-KL on a held-out WikiText test split
(Section~\ref{sec:experiments}, $-39\%$/$-33\%$); what remains is a pre-registered
three-way cost/selection/test protocol with the screened-candidate count stated.
(2)~\emph{partly done} --- a matched-cache, matched-bpp head-to-head against the
$L_1$ baseline and the $L_2$ \textsc{PrismaSCOUT} cascade is in
Table~\ref{tab:cascade} (4B); what remains is to add the closest additive priors
(Orr-style Fisher--KL \cite{orr2025}, KL-Lens-style layer KL \cite{kllens2026}),
the $L_3$ polish, and a 27B repeat, each with the wall-clock cost of producing the
allocation --- the table that lifts the novelty above synthesis.
(3)~fp32-$\dW$ with a production-cache coverage report, tying the measurement study
to the allocator. (4)~A full bpp sweep for both arms with confidence intervals, to
support or retract the ``gap widens with bits'' trend. (5)~\emph{partly done} ---
deterministic tool-call evaluation now corroborates the framing
(Section~\ref{sec:experiments}: $88$ vs.\ $84$ at matched bpp, $91$ vs.\ BF16's
$86$ for the shipped artifact); what remains is top-$k$ agreement and multi-seed
behavioral error bars. All are feasible on the $128$~GB Blackwell box with the
models already resident.

% =====================================================================
\section{Conclusion}\label{sec:conclusion}

For the $\ge 4$-bit formats that current hardware serves, mixed-precision LLM
allocation does not need explicit cross-layer interaction modeling. A single
additive per-Linear cost --- the downstream KL--Fisher contribution of the
production-rendered weight error, estimated in $O(N)$ by randomized probes ---
ranks allocations well enough that a real-KL selector finishes the job.
\textsc{Aura} realizes this: it beats the strong Fisher-weighted baseline that
was the production default at measurement time by $17.9\%$ served KL at 27B (two
disjoint corpora) and a weaker output-MSE baseline by $38\text{--}40\%$ at the 4B
knee.
The apparent additive ``overshoot'' that motivated the cross-layer apparatus is,
at these bit-widths, largely a bf16 measurement floor; the genuine residual is
small and is absorbed by selection. And the frontier's shape is rate--distortion
geometry: convex, with sharp knees living in the format menu, not the sensitivity
distribution. We have been explicit that the headline is a distributional-%
fidelity win, not a perplexity win, and that a weight-only cost is not a complete
W4A4 model --- the accounting a faithful method paper owes its reader.

% =====================================================================
\begin{thebibliography}{99}

\bibitem{prismaquant} R.~Tand. \emph{PrismaQuant: per-Linear sensitivity-driven
  mixed-precision quantization for LLMs.}
  \url{https://github.com/RobTand/prismaquant}, 2026.

\bibitem{gptq2022} E.~Frantar, S.~Ashkboos, T.~Hoefler, and D.~Alistarh. GPTQ:
  Accurate post-training quantization for generative pre-trained transformers.
  \emph{ICLR}, 2023. arXiv:2210.17323.

\bibitem{autoround2023} W.~Cheng, Y.~Lv, W.~Zhang, and H.~Shen. Optimize weight
  rounding via signed gradient descent for the quantization of LLMs.
  \emph{arXiv:2309.05516}, 2023.

\bibitem{hawqv3} Z.~Yao, Z.~Dong, Z.~Zheng, A.~Gholami, et al.
  HAWQ-V3: Dyadic Neural Network Quantization. \emph{ICML}, 2021. arXiv:2011.10680.

\bibitem{clado2023} Z.~Deng, S.~Sharify, X.~Wang, and M.~Orshansky.
  Mixed-Precision Quantization for Deep Vision Models with Integer Quadratic
  Programming (CLADO). \emph{DAC}, 2025. arXiv:2307.05657.

\bibitem{coopq} J.~Zhao, A.~Derakhshan, J.~Kana~Hyman, J.~Dong,
  S.~Abdu~Jyothi, and I.~Harris. CoopQ: Cooperative Game Inspired Layerwise
  Mixed Precision Quantization for LLMs. \emph{arXiv:2509.15455}, 2025.

\bibitem{amq} S.~Lee, S.-T.~Woo, J.~Jin, C.~Lee, and E.~Park. AMQ: Enabling
  AutoML for Mixed-precision Weight-Only Quantization of LLMs.
  \emph{arXiv:2509.12019}, 2025.

\bibitem{paroquant} Y.~Liang, H.~Chen, Z.~Zhang, S.~Han, and Z.~Liu.
  ParoQuant: Pairwise Rotation Quantization for Efficient Reasoning LLM
  Inference. \emph{ICLR}, 2026. arXiv:2511.10645.

% --- prior art that scopes Aura's novelty (added for this paper) ---
\bibitem{orr2025} D.~Orr, L.~Ribar, and C.~Luschi. Optimal Formats for Weight
  Quantisation. \emph{arXiv:2505.12988}, 2025. % closest prior art to Aura's cost

\bibitem{fimaq2025} Z.~Wu, S.~Wang, J.~Zhang, J.~Chen, and Y.~Wang.
  FIMA-Q: Post-Training Quantization for Vision Transformers by Fisher
  Information Matrix Approximation. \emph{CVPR}, 2025. arXiv:2506.11543.

\bibitem{kllens2026} J.~Kong, N.~P.~Pandey, F.~Ponzina, and T.~Rosing.
  A KL Lens on Quantization: Fast, Forward-Only Sensitivity for Mixed-Precision
  SSM-Transformer Models. \emph{arXiv:2604.13440}, 2026.

\bibitem{kneedle2011} V.~Satopaa, J.~Albrecht, D.~Irwin, and B.~Raghavan.
  Finding a ``Kneedle'' in a Haystack: Detecting Knee Points in System Behavior.
  \emph{ICDCSW}, pages 166--171, 2011.

\bibitem{ortega1998} A.~Ortega and K.~Ramchandran. Rate-Distortion Methods for
  Image and Video Compression. \emph{IEEE Signal Processing Magazine},
  15(6):23--50, 1998.

\bibitem{ocp_mx} Open Compute Project. \emph{OCP Microscaling Formats (MX)
  Specification, Version 1.0.} OCP, September 2023.

\bibitem{nvfp4_blackwell} NVIDIA Corporation. \emph{NVIDIA Blackwell
  Architecture and the NVFP4 Format.} NVIDIA Technical Briefs, 2024--2025.

\bibitem{vllm} W.~Kwon, Z.~Li, S.~Zhuang, Y.~Sheng, et al. Efficient Memory
  Management for Large Language Model Serving with PagedAttention. \emph{SOSP},
  2023.

\end{thebibliography}

\end{document}
